\ProvidesFile{tikzlibraryvc.code.tex}[2026/06/05 v1.2.4 A tikz library for 2d/3d vector computations]

% 提高数值计算精度
% https://tex.stackexchange.com/questions/521857/tikz-fpu-seems-to-be-inaccurate
\RequirePackage{xfp}
% \fpeval 默认使用弧度制, tikz 默认使用角度制, 不要使用前者作用在三角函数/反三角函数上

\usetikzlibrary{math,calc,quotes,mc}

\makeatletter

\tikzset{
  %%%%%%%%%%%%%%%%%%%%%%%%%%%%%%%%%%%%%%%%%%%%%%%
  %% Space Vector Operations
  %%%%%%%%%%%%%%%%%%%%%%%%%%%%%%%%%%%%%%%%%%%%%%%
  % exclude coorindate
  % space/exclude={\A,\B,\C,\D}
  % if A != C, result D = A
  % if B != C, result D = B
  % else throw pgferror
  space/equal/.code args={#1,#2,#3}{
    \tikzmath{
      if abs((#1{1})-(#2{1})) < 10*\EPSILON
          && abs((#1{2})-(#2{2})) < 10*\EPSILON
          && abs((#1{3})-(#2{3})) < 10*\EPSILON then {
        #3 = 1;
      } else {
        #3 = 0;
      };
    }
  },
  space/exclude/.code args={#1,#2,#3,#4}{
    \tikzset{
      space/equal={#1,#3,\cs@firstequal},
      space/equal={#2,#3,\cs@secondequal},
    }
    \tikzmath{
      if \cs@firstequal == 0 then {
        #4{1} = #1{1}; #4{2} = #1{2}; #4{3} = #1{3};
      } else {
        if \cs@secondequal == 0 then {
          #4{1} = #2{1}; #4{2} = #2{2}; #4{3} = #2{3};
        } else {
          {\pgferror{Cannot exclude coordinate}};
        };
      };
    }
  },
  % linear combination: w = a*u+b*v
  space/combine/@2 vectors/.code args={#1,#2,#3,#4,#5}{
    \tikzmath {
      #5{1} = \fpeval{(#3)*(#1{1})+(#4)*(#2{1})};
      #5{2} = \fpeval{(#3)*(#1{2})+(#4)*(#2{2})};
      #5{3} = \fpeval{(#3)*(#1{3})+(#4)*(#2{3})};
      % #5{1} = (#3)*(#1{1})+(#4)*(#2{1});
      % #5{2} = (#3)*(#1{2})+(#4)*(#2{2});
      % #5{3} = (#3)*(#1{3})+(#4)*(#2{3});
    }
  },
  % linear combination: w = a*u+b*v+c*w
  space/combine/@3 vectors/.code args={#1,#2,#3,#4,#5,#6,#7}{
    \tikzmath {
      #7{1} = \fpeval{(#4)*(#1{1})+(#5)*(#2{1})+(#6)*(#3{1})};
      #7{2} = \fpeval{(#4)*(#1{2})+(#5)*(#2{2})+(#6)*(#3{2})};
      #7{3} = \fpeval{(#4)*(#1{3})+(#5)*(#2{3})+(#6)*(#3{3})};
      % #7{1} = (#4)*(#1{1})+(#5)*(#2{1})+(#6)*(#3{1});
      % #7{2} = (#4)*(#1{2})+(#5)*(#2{2})+(#6)*(#3{2});
      % #7{3} = (#4)*(#1{3})+(#5)*(#2{3})+(#6)*(#3{3});
    }
  },
  space/combine/@argn/.initial=0,
  space/combine/.code={%dispatcher
    % compute the number of arguments
    \pgfkeyssetvalue{/tikz/space/combine/@argn}{0} % 重置计数器
    \pgfutil@for\arg:=#1\do{
      \pgfmathparse{int(add(\pgfkeysvalueof{/tikz/space/combine/@argn},1))}
      \pgfkeysalso{/tikz/space/combine/@argn = \pgfmathresult}
    }
    \ifnum \pgfkeysvalueof{/tikz/space/combine/@argn} = 5
      \tikzset{space/combine/@2 vectors= {#1}}
    \else\ifnum \pgfkeysvalueof{/tikz/space/combine/@argn} = 7
      \tikzset{space/combine/@3 vectors= {#1}}
    \else
      \pgferror{Incorrect number of arguments for 'space/combine'}
    \fi\fi
  },
  space/midpoint/.code args={#1,#2,#3}{
    \tikzset {
      space/combine={#1,#2,0.5,0.5,#3},
    }
  },
  % Vector addition: c = a+b
  space/add/origin/.code args={#1,#2,#3}{
    \tikzset {
      space/combine={#1,#2,1,1,#3},
    }
  },
  % C = P+PA+PB =-P+A+B
  space/add/.code args={#1,#2,#3,#4}{
    \tikzset {
      space/combine={#1,#2,#3,-1,1,1,#4},
    }
  },
  % Vector subtraction: c = a-b
  space/sub/origin/.code args={#1,#2,#3}{
    \tikzset {
      space/combine={#1,#2,1,-1,#3},
    }
  },
  % C = P+PA-PB = P+A-B
  space/sub/.code args={#1,#2,#3,#4}{
    \tikzset {
      space/combine={#1,#2,#3,1,1,-1,#4},
    }
  },
  % Scalar multiplication: b = k*a
  space/scale/origin/.code args={#1,#2,#3}{
    \tikzmath {
      #3{1} = \fpeval{(#2)*(#1{1})};
      #3{2} = \fpeval{(#2)*(#1{2})};
      #3{3} = \fpeval{(#2)*(#1{3})};
    }
  },
  % B = P+k (A-P) = (1-k) P+k A
  space/scale/.code args={#1,#2,#3,#4}{
    \tikzset{
      space/combine={#1,#2,(1-#3),#3,#4},
    }
  },
  % Dot product: c = a · b
  space/dot/origin/.code args={#1,#2,#3}{
    \tikzmath {
      #3 = \fpeval{#1{1}*#2{1}+#1{2}*#2{2}+#1{3}*#2{3}};
    }
  },
  % dot = PA · PB: space/dot={\P,\A,\B,\dot}
  space/dot/.code args={#1,#2,#3,#4}{
    \tikzset{
      space/sub/origin={#2,#1,\vc@pa},
      space/sub/origin={#3,#1,\vc@pb},
      space/dot/origin={\vc@pa,\vc@pb,#4},
    }
  },
  % Cross product: c = a × b
  space/cross/origin/.code args={#1,#2,#3}{
    \tikzmath{
      #3{1} = \fpeval{#1{2}*#2{3}-#1{3}*#2{2}};
      #3{2} = \fpeval{#1{3}*#2{1}-#1{1}*#2{3}};
      #3{3} = \fpeval{#1{1}*#2{2}-#1{2}*#2{1}};
    }
  },
  % C = PA x PB
  space/cross/.code args={#1,#2,#3,#4}{
    \tikzset{
      space/sub/origin={#2,#1,\vc@pa},
      space/sub/origin={#3,#1,\vc@pb},
      space/cross/origin={\vc@pa,\vc@pb,#4},
    }
  },
  % Vector length: b = |a|
  space/length/origin/.code args={#1,#2}{
    \tikzmath {
      % 无穷范数 Infinity norm
      \vc@inf = max(abs(#1{1}), abs(#1{2}), abs(#1{3}));
      % 向量长度 Vector length
      if \vc@inf < \EPSILON then {
        #2 = 0;
      } else {
        % 使用缩放法计算，防止平方溢出
        #2 = \fpeval{\vc@inf*sqrt(((#1{1})/\vc@inf)^2+((#1{2})/\vc@inf)^2+((#1{3})/\vc@inf)^2)};
      };
    }
  },
  % |PA|
  space/length/.code args={#1,#2,#3}{
    \tikzset{
      space/sub/origin={#2,#1,\vc@pa},
      space/length/origin={\vc@pa,#3},
    }
  },
  % Vector normalization: b = a/|a|
  space/normalize/origin/.code args={#1,#2}{
    \tikzset{
      space/length/origin={#1,\vc@length},
    }
    \tikzmath {
      if \vc@length > \EPSILON then {
        #2{1} = \fpeval{(#1{1})/\vc@length};
        #2{2} = \fpeval{(#1{2})/\vc@length};
        #2{3} = \fpeval{(#1{3})/\vc@length};
      } else {
        \pgferror{Zero vector detected}
        #2{1}=0; #2{2}=0; #2{3}=0;
      };
    }
  },
  % normalize: space/normalize={\P,\A,\B}
  space/normalize/.code args={#1,#2,#3}{
    \tikzset{
      space/sub/origin={#2,#1,\vc@pa},
      space/normalize/origin={\vc@pa,#3},
    }
  },
  % Scalar projection of vector 'a' on 'b' is the length (magnitude) of 
  % the component of 'a' in the direction of 'b'.
  % comp_b{a} = (a·b)/|b|
  space/component/origin/.code args={#1,#2,#3}{
    \pgfkeys{
      /tikz/space/dot/origin={#1,#2,\vc@dot},
      /tikz/space/length/origin={#2,\vc@length},
    }

    \tikzmath{
      % 投影标量
      #3 = \fpeval{\vc@dot/\vc@length};
    }
  },
  % Scalar projection of PA onto PB
  space/component/.code args={#1,#2,#3,#4}{
    \tikzset{
      space/sub/origin={#2,#1,\vc@pa},
      space/sub/origin={#3,#1,\vc@pb},
      space/component/origin={\vc@pa,\vc@pb,#4},
    }
  },
  % The vector projection of vector 'a' onto vector 'b' is the orthogonal 
  % projection of onto a straight line parallel to 'b'.
  % proj_b{a} = (a·b)/(b·b) b
  space/project/origin/.code args={#1,#2,#3}{
    \pgfkeys{
      /tikz/space/dot/origin={#1,#2,\vc@dotab},
      /tikz/space/dot/origin={#2,#2,\vc@dotbb},
      /tikz/space/scale/origin={#2,(\vc@dotab/\vc@dotbb),#3},
    }
  },
  % Vector projection of PA onto PB: PC
  % projection of A onto PB: C = P+PC
  space/project/.code args={#1,#2,#3,#4}{
    \tikzset{
      space/sub/origin={#2,#1,\vc@pa},
      space/sub/origin={#3,#1,\vc@pb},
      space/project/origin={\vc@pa,\vc@pb,\vc@pc},
      space/add/origin={#1,\vc@pc,#4},
    }
  },
  % vector rejection: rej_b{a} = a-proj_b{a}
  space/reject/origin/.code args={#1,#2,#3}{
    \tikzset{
      space/project/origin={#1,#2,\vc@proj},
      space/sub/origin={#1,\vc@proj,#3},
    }
  },
  space/reject/.code args={#1,#2,#3,#4}{
    \tikzset{
      space/sub/origin={#2,#1,\vc@pa},
      space/sub/origin={#3,#1,\vc@pb},
      space/reject/origin={\vc@pa,\vc@pb,\vc@pc},
      space/add/origin={#1,\vc@pc,#4},
    }
  },
  % vector reflection: ref_b{a} = a-2 rej_b{a} = 2 proj_b{a}-a
  space/reflect/origin/.code args={#1,#2,#3}{
    \tikzset{
      space/project/origin={#1,#2,\vc@proj},
      space/scale/origin={\vc@proj,2,\vc@temp},
      space/sub/origin={\vc@temp,#1,#3},
    }
  },
  % The symmetrical point of A with respect to line PB:
  % space/reflect={\P,\A,\B,\C}
  space/reflect/.code args={#1,#2,#3,#4}{
    \tikzset{
      space/sub/origin={#2,#1,\vc@pa},
      space/sub/origin={#3,#1,\vc@pb},
      space/reflect/origin={\vc@pa,\vc@pb,\vc@pc},
      space/add/origin={#1,\vc@pc,#4},
    }
  },
  % the angle between two vectors: angle/origin={a,b,theta}
  space/angle/origin/.code args={#1,#2,#3}{
    \tikzset{
      space/length/origin={#1,\vc@u},
      space/length/origin={#2,\vc@v},
      space/dot/origin={#1,#2,\vc@w},
    }
    
    \tikzmath{
      #3 = acos(\vc@w/(\vc@u*\vc@v));
    }
  },
  % the angle between two vectors AB and AC: angle/origin={A,B,C,theta}
  space/angle/.code args={#1,#2,#3,#4}{
    \tikzset{
      space/sub/origin={#2,#1,\vc@pa},
      space/sub/origin={#3,#1,\vc@pb},
      space/angle/origin={\vc@pa,\vc@pb,#4},
    }
  },
  % 罗德里格旋转公式 (Rodrigues' Rotation Formula)
  % 绕任意单位向量轴 $\bm{u}$ 旋转角度 $\theta$：
  % \bm{v}_{\text{rot}} = \cos\theta \bm{v} +(\bm{u} \cdot \bm{v}) (1-\cos\theta) \bm{u}+\sin\theta (\bm{u} \times \bm{v}) 
  % space/rotate/origin={u,v,theta,w}
  space/rotate/origin/.code args={#1,#2,#3,#4}{
    \tikzmath{
      % Note: \fpeval use radians, compute cos, sin immediately
      \vc@cos = cos(#3);
      \vc@sin = sin(#3);
    }
    
    \tikzset{
      space/normalize/origin={#1,\vc@U},
      space/cross/origin={\vc@U,#2,\vc@cross},
      space/dot/origin={\vc@U,#2,\vc@dot},
      space/combine={#2,\vc@cross,\vc@U,
        \vc@cos,
        \vc@sin,
        \vc@dot*(1-\vc@cos),
        #4},
    }
    \message{U:(\vc@U{1},\vc@U{2},\vc@U{3})^^J}
    \message{V:(#2{1},#2{2},#2{3})^^J}
    \message{UxV:(\vc@cross{1},\vc@cross{2},\vc@cross{3})^^J}
    \message{U·V:\vc@dot^^J}
  },
  % PV rotates around PU by an angle theta to obtain PW
  % space/rotate={\P,\U,\V,\ang,\W}
  space/rotate/.code args={#1,#2,#3,#4,#5}{
    \tikzset{
      space/sub/origin={#2,#1,\vc@Ua},
      space/sub/origin={#3,#1,\vc@Va},
      space/rotate/origin={\vc@Ua,\vc@Va,#4,\vc@W},
      space/add/origin={#1,\vc@W,#5},
    }
  },
  % Gram-Schmidt Orthogonalization
  % The Gram-Schmidt orthogonalization process is a mathematical 
  % algorithm used to transform a set of linearly independent 
  % vectors into an orthogonal (or orthonormal) basis.
  % space/orthonormal basis={\U,\V,\W},
  % \U,\V,\W: unit vectors, \U is the input, \V,\W are the outputs
  space/orthonormal basis/.code args={#1,#2,#3}{
    \tikzmath{
      % 为避免 向量 a 与向量 u 接近平行导致数值误差,
      % 采取下面的策略确保向量 a 与向量 u 夹角不太小, 提高数值稳定性
      if abs(#1{1}) <= abs(#1{2}) && abs(#1{1}) <= abs(#1{3}) then {
        \vc@a{1} = 1; \vc@a{2} = 0; \vc@a{3} = 0;
      } else {
        if abs(#1{2}) <= abs(#1{3}) then {
          \vc@a{1} = 0; \vc@a{2} = 1; \vc@a{3} = 0;    
        } else {
          \vc@a{1} = 0; \vc@a{2} = 0; \vc@a{3} = 1;
        };
      };
    }
    \tikzset{
      % Gram-Schmidt Orthogonalization Process
      space/reject/origin={\vc@a,#1,\vc@v},
      space/normalize/origin={\vc@v,#2},
      space/cross/origin={#1,#2,#3},
    }
  },
  % Intersection of a plane and a line
  % space/intersect pl={\P,\n,\A,\B,\C}
  % \P - a point on the plane
  % \n - the normal vector of the plane
  % \A, \B - two points on the line
  % \C - the intersection
  % raise error if the line is parallel to a plane or lies on the plane
  space/@assert_not_parallel/.code args={#1,#2,#3,#4}{
    \tikzset{
      space/normalize/origin={#2,\vc@n},
      space/normalize={#3,#4,\vc@l},
      space/dot/origin={\vc@n,\vc@l,\vc@dot},
    }

    \tikzmath{
      if abs(\vc@dot) < \EPSILON then {
        {\pgferror{The line is parallel to a plane or lies on the plane}};
      };
    }
  },
  space/intersect pl/.code args={#1,#2,#3,#4,#5}{
    \tikzset{
      space/@assert_not_parallel={#1,#2,#3,#4},
      space/sub/origin={#4,#3,\vc@ab},
      space/sub/origin={#1,#3,\vc@ap},
      space/dot/origin={\vc@n,\vc@ab,\vc@dota},
      space/dot/origin={\vc@n,\vc@ap,\vc@dotb},
      space/scale={#3,#4,\fpeval{\vc@dotb/\vc@dota},#5},
    }
  },
  % perspective transformation matrix or homography matrix
  % space/perspective transform={\Center,\S,\m,\T,\n,\H}
  % S - the origin of source plane, 源平面
  % m - the normal vector of source plane, 源平面
  % T - the origin of target plane, 目标平面
  % n - the normal vector of target plane, 目标平面
  space/perspective transform/.code args={#1,#2,#3,#4,#5,#6}{
    \tikzset{
      space/normalize/origin={#3,\vc@m},
      space/normalize/origin={#5,\vc@n},
      space/sub/origin={#2,#1,\vc@SC},
      space/sub/origin={#4,#1,\vc@TC},
      space/orthonormal basis={\vc@m,\vc@SU,\vc@SV},
      space/orthonormal basis={\vc@n,\vc@TU,\vc@TV},
      space/dot/origin={\vc@n,\vc@SU,\vc@Aa},
      space/dot/origin={\vc@n,\vc@SV,\vc@Ab},
      space/dot/origin={\vc@n,\vc@SC,\vc@Ac},
      space/dot/origin={\vc@n,\vc@TC,\vc@Ad},
      space/dot/origin={\vc@TU,\vc@SU,\vc@Ba},
      space/dot/origin={\vc@TU,\vc@SV,\vc@Bb},
      space/dot/origin={\vc@TU,\vc@SC,\vc@Bc},
      space/dot/origin={\vc@TU,\vc@TC,\vc@Bd},
      space/dot/origin={\vc@TV,\vc@SU,\vc@Ca},
      space/dot/origin={\vc@TV,\vc@SV,\vc@Cb},
      space/dot/origin={\vc@TV,\vc@SC,\vc@Cc},
      space/dot/origin={\vc@TV,\vc@TC,\vc@Cd},  
    }
    \tikzmath{
      #6{1,1} = \fpeval{\vc@Ad*\vc@Ba-\vc@Aa*\vc@Bd};
      #6{1,2} = \fpeval{\vc@Ad*\vc@Bb-\vc@Ab*\vc@Bd};
      #6{1,3} = \fpeval{\vc@Ad*\vc@Bc-\vc@Ac*\vc@Bd};
      #6{2,1} = \fpeval{\vc@Ad*\vc@Ca-\vc@Aa*\vc@Cd};
      #6{2,2} = \fpeval{\vc@Ad*\vc@Cb-\vc@Ab*\vc@Cd};
      #6{2,3} = \fpeval{\vc@Ad*\vc@Cc-\vc@Ac*\vc@Cd};
      #6{3,1} = \vc@Aa;
      #6{3,2} = \vc@Ab;
      #6{3,3} = \vc@Ac;
    }
  },
  %%%%%%%%%%%%%%%%%%%%%%%%%%%%%%%%%%%%%%%%%%%%%%%
  %% Plane Vector Operations
  %%%%%%%%%%%%%%%%%%%%%%%%%%%%%%%%%%%%%%%%%%%%%%%
  % exclude coorindate
  % plane/exclude={\A,\B,\C,\D}
  % if A != C, result D = A
  % if B != C, result D = B
  % else throw pgferror
  plane/equal/.code args={#1,#2,#3}{
    \tikzmath{
      if abs((#1{1})-(#2{1})) < 10*\EPSILON
          && abs((#1{2})-(#2{2})) < 10*\EPSILON then {
        #3 = 1;
      } else {
        #3 = 0;
      };
    }
  },
  plane/exclude/.code args={#1,#2,#3,#4}{
    \tikzset{
      plane/equal={#1,#3,\cs@firstequal},
      plane/equal={#2,#3,\cs@secondequal},
    }
    \tikzmath{
      if \cs@firstequal == 0 then {
        #4{1} = #1{1}; #4{2} = #1{2};
      } else {
        if \cs@secondequal == 0 then {
          #4{1} = #2{1}; #4{2} = #2{2};
        } else {
          {\pgferror{Cannot exclude coordinate}};
        };
      };
    }
    % \message{A: (#1{1},#1{2})^^J}
    % \message{B: (#2{1},#2{2})^^J}
    % \message{C: (#3{1},#3{2})^^J}
    % \message{ans: (#4{1},#4{2})^^J}
  },
  % linear combination: w = a*u+b*v
  plane/combine/@2 vectors/.code args={#1,#2,#3,#4,#5}{
    % \message{===== Linear Combination =====^^J}
    % \message{A: (#1{1},#1{2}) a: #3 ^^J}
    % \message{B: (#2{1},#2{2}) b: #4 ^^J}

    \tikzmath {
      #5{1} = \fpeval{(#3)*(#1{1})+(#4)*(#2{1})};
      #5{2} = \fpeval{(#3)*(#1{2})+(#4)*(#2{2})};
    }

    % \message{ans: (#5{1},#5{2})^^J}
  },
  % linear combination: w = a*u+b*v+c*w
  plane/combine/@3 vectors/.code args={#1,#2,#3,#4,#5,#6,#7}{
    \tikzmath {
      #7{1} = \fpeval{(#4)*(#1{1})+(#5)*(#2{1})+(#6)*(#3{1})};
      #7{2} = \fpeval{(#4)*(#1{2})+(#5)*(#2{2})+(#6)*(#3{2})};
    }
  },
  plane/combine/@argn/.initial=0,
  plane/combine/.code={%dispatcher
    % compute the number of arguments
    \pgfkeyssetvalue{/tikz/plane/combine/@argn}{0} % 重置计数器
    \pgfutil@for\arg:=#1\do{
      \pgfmathparse{int(add(\pgfkeysvalueof{/tikz/plane/combine/@argn},1))}
      \pgfkeysalso{/tikz/plane/combine/@argn = \pgfmathresult}
    }
    \ifnum \pgfkeysvalueof{/tikz/plane/combine/@argn} = 5
      \tikzset{plane/combine/@2 vectors= {#1}}
    \else\ifnum \pgfkeysvalueof{/tikz/plane/combine/@argn} = 7
      \tikzset{plane/combine/@3 vectors= {#1}}
    \else
      \pgferror{Incorrect number of arguments for 'plane/combine'}
    \fi\fi
  },
  plane/midpoint/.code args={#1,#2,#3}{
    \tikzset {
      plane/combine={#1,#2,0.5,0.5,#3},
    }
  },
  % Vector addition: c = a+b
  plane/add/origin/.code args={#1,#2,#3}{
    \tikzset {
      plane/combine={#1,#2,1,1,#3},
    }

    % \typeout{=========================}
    % \typeout{vector a: (#1{1},#1{2})}
    % \typeout{vector b: (#2{1},#2{2})}
    % \typeout{vector a+b: (#3{1},#3{2})}
    % \typeout{=========================}
  },
  % C = P+PA+PB = -P+A+B
  plane/add/.code args={#1,#2,#3,#4}{
    \tikzset {
      plane/combine={#1,#2,#3,-1,1,1,#4},
    }
  },
  % Vector subtraction: c = a-b
  plane/sub/origin/.code args={#1,#2,#3}{
    \tikzset {
      plane/combine={#1,#2,1,-1,#3},
    }
  },
  % C = P+PA-PB = P+A-B
  plane/sub/.code args={#1,#2,#3,#4}{
    \tikzset {
      plane/combine={#1,#2,#3,1,1,-1,#4},
    }
  },
  % Scalar multiplication: b = k*a
  plane/scale/origin/.code args={#1,#2,#3}{
    \tikzmath {
      #3{1} = \fpeval{(#2)*(#1{1})};
      #3{2} = \fpeval{(#2)*(#1{2})};
    }
  },
  % B = P+k (A-P) = (1-k) P+k A
  plane/scale/.code args={#1,#2,#3,#4}{
    \tikzset{
      plane/combine={#1,#2,(1-#3),#3,#4},
    }
  },
  % Rotation: rotate/origin={a,angle,b}
  plane/rotate/origin/.code args={#1,#2,#3}{
    \tikzmath {
      #3{1} = #1{1}*cos(#2)-#1{2}*sin(#2);
      #3{2} = #1{1}*sin(#2)+#1{2}*cos(#2);
    }
  },
  % vector rotation: plane/rotate={P,A,angle,B} 
  plane/rotate/.code args={#1,#2,#3,#4}{
    \tikzset{
      plane/sub/origin={#2,#1,\vc@pa},
      plane/rotate/origin={\vc@pa,#3,\vc@pb},
      plane/add/origin={#1,\vc@pb,#4},
    }
  },
  % Dot product: c = a · b
  plane/dot/origin/.code args={#1,#2,#3}{
    \tikzmath {
      #3 = \fpeval{#1{1}*#2{1}+#1{2}*#2{2}};
    }
  },
  % dot = PA · PB: plane/dot={\P,\A,\B,\dot}
  plane/dot/.code args={#1,#2,#3,#4}{
    \tikzset{
      plane/sub/origin={#2,#1,\vc@pa},
      plane/sub/origin={#3,#1,\vc@pb},
      plane/dot/origin={\vc@pa,\vc@pb,#4},
    }
  },
  % Vector length: b = |a|
  plane/length/origin/.code args={#1,#2}{
    \tikzmath {
      % 无穷范数 Infinity norm
      \vc@inf = max(abs(#1{1}), abs(#1{2}));
      % 向量长度 Vector length
      if \vc@inf < \EPSILON then {
        #2 = 0;
      } else {
        % 使用缩放法计算，防止平方溢出
        #2 = \fpeval{\vc@inf*sqrt(((#1{1})/\vc@inf)^2+((#1{2})/\vc@inf)^2)};
      };
    }
  },
  plane/length/@segment/.code args={#1,#2,#3}{
    \tikzset{
      plane/sub/origin={#2,#1,\vc@pa},
      plane/length/origin={\vc@pa,#3},
    }
  },
  plane/length/@triangle/.code args={#1,#2,#3,#4,#5,#6}{
    \tikzset{
      plane/length/@segment={#2,#3,#4},
      plane/length/@segment={#3,#1,#5},
      plane/length/@segment={#1,#2,#6},
    }
    
    \tikzmath{
      % {\message{Triangle sides: a=\vc@a,b=\vc@b,c=\vc@c^^J}};
      if (\vc@a+\vc@b > \vc@c) && (\vc@b+\vc@c > \vc@a) && (\vc@c+\vc@a > \vc@b) then {
        %pass
      } else {
        {\pgferror{Three points are collinear}};
      };
    }
  },
  plane/length/@argn/.initial=0,
  plane/length/.code={%dispatcher
    % compute the number of arguments
    \pgfkeyssetvalue{/tikz/plane/length/@argn}{0} % 重置计数器
    \pgfutil@for\arg:=#1\do{
      \pgfmathparse{int(add(\pgfkeysvalueof{/tikz/plane/length/@argn},1))}
      \pgfkeysalso{/tikz/plane/length/@argn = \pgfmathresult}
    }
    \ifnum \pgfkeysvalueof{/tikz/plane/length/@argn} = 3
      \tikzset{plane/length/@segment = {#1}}
    \else\ifnum \pgfkeysvalueof{/tikz/plane/length/@argn} = 6
      \tikzset{plane/length/@triangle = {#1}}
    \else
      \pgferror{Incorrect number of arguments for 'plane/length'}
    \fi\fi
  },
  % Vector normalization: b = a/|a|
  plane/normalize/origin/.code args={#1,#2}{
    \tikzset{
      plane/length/origin={#1,\vc@length},
    }
    \tikzmath {
      if \vc@length > \EPSILON then {
        #2{1} = \fpeval{(#1{1})/\vc@length};
        #2{2} = \fpeval{(#1{2})/\vc@length};
      } else {
        #2{1}=0; #2{2}=0;
        \pgferror{Zero vector detected}
      };
    }
  },
  % normalize: plane/normalize={\P,\A,\B}
  plane/normalize/.code args={#1,#2,#3}{
    \tikzset{
      plane/sub/origin={#2,#1,\vc@pa},
      plane/normalize/origin={\vc@pa,#3},
    }
  },
  % Scalar projection of vector 'a' on 'b' is the length (magnitude) of 
  % the component of 'a' in the direction of 'b'.
  % comp_b{a} = (a·b)/|b|
  plane/component/origin/.code args={#1,#2,#3}{
    \pgfkeys{
      /tikz/plane/dot/origin={#1,#2,\vc@dot},
      /tikz/plane/length/origin={#2,\vc@length},
    }

    \tikzmath{
      % 投影标量
      #3 = \fpeval{\vc@dot/\vc@length};
    }
  },
  % Scalar projection of PA onto PB
  plane/component/.code args={#1,#2,#3,#4}{
    \tikzset{
      plane/sub/origin={#2,#1,\vc@pa},
      plane/sub/origin={#3,#1,\vc@pb},
      plane/component/origin={\vc@pa,\vc@pb,#4},
    }
  },
  % The vector projection of vector 'a' onto vector 'b' is the orthogonal 
  % projection of onto a straight line parallel to 'b'.
  % proj_b{a} = (a·b)/(b·b) b
  plane/project/origin/.code args={#1,#2,#3}{
    \pgfkeys{
      /tikz/plane/dot/origin={#1,#2,\vc@dotab},
      /tikz/plane/dot/origin={#2,#2,\vc@dotbb},
      /tikz/plane/scale/origin={#2,(\vc@dotab/\vc@dotbb),#3},
    }
  },
  % Vector projection of PA onto PB: PC
  % projection of A onto PB: C = P+PC
  plane/project/.code args={#1,#2,#3,#4}{
    \tikzset{
      plane/sub/origin={#2,#1,\vc@pa},
      plane/sub/origin={#3,#1,\vc@pb},
      plane/project/origin={\vc@pa,\vc@pb,\vc@pc},
      plane/add/origin={#1,\vc@pc,#4},
    }
  },
  % vector rejection: rej_b{a} = a-proj_b{a}
  plane/reject/origin/.code args={#1,#2,#3}{
    \tikzset{
      plane/project/origin={#1,#2,\vc@proj},
      plane/sub/origin={#1,\vc@proj,#3},
    }
  },
  plane/reject/.code args={#1,#2,#3,#4}{
    \tikzset{
      plane/sub/origin={#2,#1,\vc@pa},
      plane/sub/origin={#3,#1,\vc@pb},
      plane/reject/origin={\vc@pa,\vc@pb,\vc@pc},
      plane/add/origin={#1,\vc@pc,#4},
    }
  },
  % vector reflection: ref_b{a} = a-2 rej_b{a} = 2 proj_b{a}-a
  plane/reflect/origin/.code args={#1,#2,#3}{
    \tikzset{
      plane/project/origin={#1,#2,\vc@proj},
      plane/scale/origin={\vc@proj,2,\vc@temp},
      plane/sub/origin={\vc@temp,#1,#3},
    }
  },
  % The symmetrical point of A with respect to line PB:
  % space/reflect={\P,\A,\B,\C}
  plane/reflect/.code args={#1,#2,#3,#4}{
    \tikzset{
      plane/sub/origin={#2,#1,\vc@pa},
      plane/sub/origin={#3,#1,\vc@pb},
      plane/reflect/origin={\vc@pa,\vc@pb,\vc@pc},
      plane/add/origin={#1,\vc@pc,#4},
    }
  },
  % the oriented (or signed) angle between two vectors: angle/origin={a,b,theta}
  plane/angle/origin/.code args={#1,#2,#3}{
    \tikzmath{
      % DO NOT USE \fpeval, which returns angle in radians
      #3 = atan2((#1{1})*(#2{2})-(#1{2})*(#2{1}), (#1{1})*(#2{1})+(#1{2})*(#2{2}));
    }
  },
  % the oriented (or signed) angle between two vectors AB and AC: angle/origin={A,B,C,theta}
  plane/angle/.code args={#1,#2,#3,#4}{
    \tikzset{
      plane/sub/origin={#2,#1,\vc@pa},
      plane/sub/origin={#3,#1,\vc@pb},
      plane/angle/origin={\vc@pa,\vc@pb,#4},
    }
  },
  %%%%%%%%%%%%%%%%%%%%%%%%%%%%%%%%%%%%%%%%%%%%%%%
  %% Triangle
  %%%%%%%%%%%%%%%%%%%%%%%%%%%%%%%%%%%%%%%%%%%%%%%
  % perimeter
  % plane/perimeter={\A,\B,\C,\p}
  plane/perimeter/.code args={#1,#2,#3,#4}{
    \tikzset{
      plane/length={#1,#2,#3,\vc@a,\vc@b,\vc@c},
    }

    \tikzmath{
      #4 = \fpeval{\vc@a+\vc@b+\vc@c};
    }
  },
  % Heron's Formula
  plane/area/.code args={#1,#2,#3,#4}{
    \tikzset{
      plane/length={#1,#2,#3,\vc@a,\vc@b,\vc@c},
    }

    \tikzmath{
      \vc@s = \fpeval{0.5*(\vc@a+\vc@b+\vc@c)};
      \vc@ka = \fpeval{\vc@a/\vc@s};
      \vc@kb = \fpeval{\vc@b/\vc@s};
      \vc@kc = \fpeval{\vc@c/\vc@s};
      #4 = \fpeval{(\vc@s)^2*sqrt((1-\vc@ka)*(1-\vc@kb)*(1-\vc@kc))};
    }
  },
  plane/circumradius/.code args={#1,#2,#3,#4}{
    \tikzset{
      % byproducts: sides(\vc@a,\vc@b,\vc@c) and semiperimeter(\vc@s)
      plane/area={#1,#2,#3,\vc@area},
    }

    \tikzmath{
      #4 = \fpeval{(0.25/\vc@area)*\vc@a*\vc@b*\vc@c};
    }
  },
  plane/inradius/.code args={#1,#2,#3,#4}{
    \tikzset{
      % byproducts: sides(\vc@a,\vc@b,\vc@c) and semiperimeter(\vc@s)
      plane/area={#1,#2,#3,\vc@area},
    }

    \message{a,b,c:\vc@a,\vc@b,\vc@c^^J};
    \message{semiperimeter:\vc@s^^J};
    \message{area:\vc@area^^J};

    \tikzmath{
      #4 = \fpeval{\vc@area/\vc@s};
    }
  },
  plane/exradius/.code args={#1,#2,#3,#4}{
    \tikzset{
      % byproducts: sides(\vc@a,\vc@b,\vc@c) and semiperimeter(\vc@s)
      plane/area={#1,#2,#3,\vc@area},
    }

    \tikzmath{
      #4 = \fpeval{\vc@area/(\vc@s-\vc@a)};
    }
  },
  %%%%%%%%%%%%%%%%%%%%%%%%%%%%%%%%%%%%%%%%%%%%%%%
  %% Triangle Centers
  %%%%%%%%%%%%%%%%%%%%%%%%%%%%%%%%%%%%%%%%%%%%%%%
  % barycentric coordinatess
  % P = (x A+y B+z C)/(x+y+z)
  % plane/barycenter={\A,\B,\C,\x,\y,\z,\P}
  % a = x/(x+y+z), b = y/(x+y+z), c = z/(x+y+z)
  % P = a A+b B+c C
  plane/barycenter/.code args={#1,#2,#3,#4,#5,#6,#7}{
    \tikzmath{
      \vc@sum = \fpeval{(#4)+(#5)+(#6)};
      if abs(\vc@sum) < \EPSILON then {
        {\pgferror{Illeagal ratios in 'plane/barycenter'}};
      };
    }

    \tikzset{
      plane/combine={#1,#2,#3,(#4/\vc@sum),(#5/\vc@sum),(#6/\vc@sum),#7},
    }
  },
  % incenter: plane/incenter={\A,\B,\C,\I}
  plane/incenter/.code args={#1,#2,#3,#4}{
    \tikzset{
      plane/length={#1,#2,#3,\vc@a,\vc@b,\vc@c},
    }

    \tikzset{
      plane/barycenter={#1,#2,#3,
        (\vc@a),
        (\vc@b),
        (\vc@c),#4},
    }
  },
  % excenter: plane/excenter={\A,\B,\C,\Ja}, returns excenter opposite to the vertex A
  plane/excenter/.code args = {#1,#2,#3,#4}{ 
    \tikzset{
      plane/length={#1,#2,#3,\vc@a,\vc@b,\vc@c},
    }

    \tikzset{
      plane/barycenter={#1,#2,#3,
        (-\vc@a),
        (\vc@b),
        (\vc@c),#4},
    }
  },
  % plane/circumcenter = {\A,\B,\C,\O}
  plane/circumcenter/.code args = {#1,#2,#3,#4}{ 
    \tikzset{
      plane/length={#1,#2,#3,\vc@a,\vc@b,\vc@c},
    }

    \tikzmath{
      \vc@m = max(\vc@a,\vc@b,\vc@c);
      \vc@a = \fpeval{\vc@a/\vc@m};
      \vc@b = \fpeval{\vc@b/\vc@m};
      \vc@c = \fpeval{\vc@c/\vc@m};
      \vc@temp{1} = \fpeval{\vc@a^2*(\vc@b^2+\vc@c^2-\vc@a^2)};
      \vc@temp{2} = \fpeval{\vc@b^2*(\vc@c^2+\vc@a^2-\vc@b^2)};
      \vc@temp{3} = \fpeval{\vc@c^2*(\vc@a^2+\vc@b^2-\vc@c^2)};
      % {\message{Barycentric coordinates: (\vc@temp{1},\vc@temp{2},\vc@temp{3})^^J}};
    }

    \tikzset{
      plane/barycenter={#1,#2,#3,
        (\vc@temp{1}),
        (\vc@temp{2}),
        (\vc@temp{3}),#4},
    }
  },
  % plane/orthocenter={\A,\B,\C,\H}
  plane/orthocenter/.code args = {#1,#2,#3,#4}{ 
    \tikzset{
      plane/length={#1,#2,#3,\vc@a,\vc@b,\vc@c},
    }

    \tikzmath{
      \vc@m = max(\vc@a,\vc@b,\vc@c);
      \vc@a = \fpeval{\vc@a/\vc@m};
      \vc@b = \fpeval{\vc@b/\vc@m};
      \vc@c = \fpeval{\vc@c/\vc@m};
      \vc@t{1} = \fpeval{(\vc@a^2+\vc@b^2-\vc@c^2)};
      \vc@t{2} = \fpeval{(\vc@b^2+\vc@c^2-\vc@a^2)};
      \vc@t{3} = \fpeval{(\vc@c^2+\vc@a^2-\vc@b^2)};
      \vc@temp{1} = \fpeval{\vc@t{1}*\vc@t{3}};
      \vc@temp{2} = \fpeval{\vc@t{2}*\vc@t{1}};
      \vc@temp{3} = \fpeval{\vc@t{3}*\vc@t{2}};
    }

    \tikzset{
      plane/barycenter={#1,#2,#3,
        (\vc@temp{1}),
        (\vc@temp{2}),
        (\vc@temp{3}),#4},
    }

    % \message{Cartesian coordinates: (#4{1},#4{2})^^J}
  },
  % plane/centroid={\A,\B,\C,\G}
  plane/centroid/.code args={#1,#2,#3,#4}{
    \tikzset{
      plane/barycenter={#1,#2,#3,
        1,1,1,#4},
    }
  },
  %%%%%%%%%%%%%%%%%%%%%%%%%%%%%%%%%%%%%%%%%%%%%%%
  %% Lines and Circles
  %%%%%%%%%%%%%%%%%%%%%%%%%%%%%%%%%%%%%%%%%%%%%%%
  % projection of origin onto the line
  % plane/anchor={\a,\b,\c,\P}
  plane/anchor/.code args={#1,#2,#3,#4}{
    \tikzmath{
      % 原点在直线上的垂足
      \vc@square = \fpeval{(#1)^2+(#2)^2};
      #4{1} = \fpeval{-(#1)*(#3)/\vc@square};
      #4{2} = \fpeval{-(#2)*(#3)/\vc@square};
    }
  },
  % direction of line:
  % plane/direct={\a,\b,\c,\Direct}
  plane/direct/.code args={#1,#2,#3,#4}{
    \tikzmath{
      if #1 > 0 then {
        #4{1} = -(#2);
        #4{2} = #1;
      } else {
        if abs(#1) > \EPSILON then {
          #4{1} = #2;
          #4{2} = -(#1);
        } else {
          #4{1} = abs(#2);
          #4{2} = 0;
        };
      };
    }
  },
  % intersection of line AB and line CD
  % plane/intersect ll={\A,\B,\C,\D,\P}
  plane/intersect ll/.code args={#1,#2,#3,#4,#5}{
    \tikzmath{
      % (\yb-\ya) x-(\xb-\xa) y = \xa*(\yb-\ya)-\ya*(\xb-\xa)
      % (\yd-\yc) x-(\xd-\xc) y = \xc*(\yd-\yc)-\yc*(\xd-\xc)
      \vc@a{1,1} = \fpeval{#2{2}-#1{2}}; \vc@a{1,2} = \fpeval{-(#2{1}-#1{1})}; 
      \vc@b{1,1} = \fpeval{#1{1}*\vc@a{1,1}+#1{2}*\vc@a{1,2}};
      \vc@a{2,1} = \fpeval{#4{2}-#3{2}}; \vc@a{2,2} = \fpeval{-(#4{1}-#3{1})}; 
      \vc@b{2,1} = \fpeval{#3{1}*\vc@a{2,1}+#3{2}*\vc@a{2,2}};
    }
    \tikzset{
      mat/solve={\vc@a,\vc@b,2,1,\vc@ans},
    }
    \tikzmath{
      #5{1} = \vc@ans{1,1};
      #5{2} = \vc@ans{2,1};
    }
    % \typeout{=========================}
    % \typeout{Line-Line Intersection:}
    % \tikzset{
    %   mat/stdout={\vc@a,2,2},
    %   mat/stdout={\vc@b,2,1},
    %   mat/stdout={\vc@ans,2,1},
    % }
    % \typeout{=========================}
    % \coordinate (ll) at (\c{1,1},\c{2,1});
  },
  % intersection of circle and line (through point A and B)
  % plane/intersect cl={\O,\r,\A,\B,\Pa,\Pb}
  plane/intersect cl/.code args={#1,#2,#3,#4,#5,#6}{
    \tikzset{
      plane/normalize={#3,#4,\vc@u},
      % projection of circle center onto line
      plane/project={#3,#1,#4,\vc@mid},
      plane/length={#1,\vc@mid,\vc@d},
    }
    \tikzmath{
      if \vc@d > abs(#2) then {
        % 在 tikzmath 中用 {} 将 \pgferror 包裹起来，将其作为普通 TeX 代码块执行
        {\pgferror{The circle and the line do not intersect}};
      };
      \vc@l = \fpeval{sqrt((#2)^2-(\vc@d)^2)};
    }
    \tikzset{
      plane/scale/origin={\vc@u,\vc@l,\vc@v},
      plane/sub/origin={\vc@mid,\vc@v,#5},
      plane/add/origin={\vc@mid,\vc@v,#6},
    }
  },
    % intersection of circle and line (a x+b y+c = 0)
  % plane/intersect cn={\O,\r,\a,\b,\c,\Pa,\Pb}
  plane/intersect cn/.code args={#1,#2,#3,#4,#5,#6,#7}{
    \tikzset{
      plane/anchor={#3,#4,#5,\vc@P},
      plane/direct={#3,#4,#5,\vc@direct},
      plane/add/origin={\vc@P,\vc@direct,\vc@Q},
      plane/intersect cl={#1,#2,\vc@P,\vc@Q,#6,#7},
    }
  },
  % The radical axis of two non-concentric circles is a straight line 
  % perpendicular to the line of centers (the line connecting the centers 
  % of the circles). 
  % plane/radical axis={\Oa,\ra,\Ob,\rb,\H,\U}
  % \H: intersection of radical axis and line of centers
  % H = O1+1/2 (d^2+r1^2-r2^2)/d^2 (O2-O1), d = |O1O2|
  % \U: unit direction vector
  plane/radical axis/@point-direction/.code args={#1,#2,#3,#4,#5,#6}{
    \tikzset{
      plane/length={#1,#3,\vc@c},
    }

    % determine if there is a radical axis
    \tikzmath{
      \vc@a = abs(#2);
      \vc@b = abs(#4);
      if \vc@c < \EPSILON then {
        {\pgferror{Concentric circles have no radical axis}};
      };
      \vc@scale = \fpeval{1/2*(\vc@c^2+\vc@a^2-\vc@b^2)/\vc@c^2};
    }

    \tikzset{
      plane/scale={#1,#3,\vc@scale,#5},
      plane/normalize={#1,#3,\vc@n},
      plane/rotate/origin={\vc@n,90,#6},
    }

    \tikzmath{
      if abs(#6{2}) > \EPSILON then {
        if #6{2} < 0 then {
          #6{1} = -(#6{1});
          #6{2} = -(#6{2});
        };
      } else {
        if #6{1} < 0 then {
          #6{1} = -(#6{1});
          #6{2} = 0;
        };
      };
    }
  },
  % plane/radical axis={\Oa,\ra,\Ob,\rb,\a,\b,\c}
  plane/radical axis/@equation/.code args={#1,#2,#3,#4,#5,#6,#7}{
    \tikzmath{
      #5 = \fpeval{2*(#3{1}-#1{1})};
      #6 = \fpeval{2*(#3{2}-#1{2})};
      #7 = \fpeval{((#1{1})^2+(#1{2})^2-(#2)^2)-((#3{1})^2+(#3{2})^2-(#4)^2)};
    }
  },
  plane/radical axis/@argn/.initial=0,
  plane/radical axis/.code={%dispatcher
    % compute the number of arguments
    \pgfkeyssetvalue{/tikz/plane/radical axis/@argn}{0} % 重置计数器
    \pgfutil@for\arg:=#1\do{
      \pgfmathparse{int(add(\pgfkeysvalueof{/tikz/plane/radical axis/@argn},1))}
      \pgfkeysalso{/tikz/plane/radical axis/@argn = \pgfmathresult}
    }
    \ifnum \pgfkeysvalueof{/tikz/plane/radical axis/@argn} = 6
      \tikzset{plane/radical axis/@point-direction = {#1}}
    \else\ifnum \pgfkeysvalueof{/tikz/plane/radical axis/@argn} = 7
      \tikzset{plane/radical axis/@equation = {#1}}
    \else
      \pgferror{Incorrect number of arguments for 'plane/radical axis'}
    \fi\fi
  },
  % intersection of circle and circle
  % plane/intersect cc={\Oa,\ra,\Ob,\rb,\Pa,\Pb}
  plane/intersect cc/.code args={#1,#2,#3,#4,#5,#6}{
    \tikzset{
      plane/radical axis={#1,#2,#3,#4,\vc@P,\vc@U},
      plane/add/origin={\vc@P,\vc@U,\vc@Q},
      plane/intersect cl={#1,#2,\vc@P,\vc@Q,#5,#6},
    }
  },
  % tangent lines from an external point to a circle
  % return two tangent points: plane/tangents={\O,\r,\P,\Ta,\Tb}
  plane/tangents/.code args={#1,#2,#3,#4,#5}{
    \tikzset{
      plane/length={#1,#3,\vc@d},
      plane/scale={#1,#3,0.5,\vc@center},
      plane/intersect cc={#1,#2,\vc@center,(\vc@d/2),#4,#5},
    }
  },
  % inverse point
  % plane/inverse={\O,\r,\P,\Q}
  plane/inverse/.code args={#1,#2,#3,#4}{
    \tikzset{
      plane/length={#1,#3,\vc@dist},
      plane/scale={#1,#3,\fpeval{(#2)^2/(\vc@dist)^2},#4},
    }
  },
  % External Center of Similitude
  % plane/external center={\Oa,\ra,\Ob,\rb,\center}
  plane/external center/.code args={#1,#2,#3,#4,#5}{
    \tikzset{
      plane/scale={#1,#3,(#2/(#2-#4)),#5},
    }
  },
  % Internal Center of Similitude
  % plane/internal center={\Oa,\ra,\Ob,\rb,\center}
  plane/internal center/.code args={#1,#2,#3,#4,#5}{
    \tikzset{
      plane/scale={#1,#3,(#2/(#2+#4)),#5},
    }
  },
}

\makeatother
